\documentclass[border=6pt]{standalone}
\usepackage{tikz}
\usepackage{amsmath}
\usepackage{xeCJK}
\setCJKmainfont{Noto Sans CJK SC}
\usetikzlibrary{arrows.meta,positioning,calc}
\begin{document}
\begin{tikzpicture}[
  font=\small,
  step/.style={draw,very thick,minimum width=3.6cm,minimum height=1.0cm,align=center},
  init/.style={draw,double,very thick,minimum width=3.6cm,minimum height=1.0cm,align=center},
  act/.style={draw,minimum width=3.4cm,minimum height=1.0cm,align=left,font=\footnotesize},
  every node/.style={},
  ]
  % steps on the left column
  \node[init] (s0) at (0,10) {S0\quad 初始步（停机）};
  \node[step] (s1) at (0,7.6) {S1\quad 步1};
  \node[step] (s2) at (0,5.2) {S2\quad 步2};
  \node[step] (s3) at (0,2.8) {S3\quad 步3};
  \node[step] (s4) at (0,0.4) {S4\quad 步4（暂停）};

  % action boxes on the right of each step
  \node[act,right=1.4cm of s1] (a1) {P1(Q4.0)=1\\ A组(Q5.0$\sim$Q5.6)=1};
  \node[act,right=1.4cm of s2] (a2) {P2(Q4.1)=1\\ B组(Q6.0$\sim$Q6.6)=1};
  \node[act,right=1.4cm of s3] (a3) {P1,P2=1\\ A组,B组=1};
  \node[act,right=1.4cm of s4] (a4) {全部输出=0};
  \draw (s1.east)--(a1.west);
  \draw (s2.east)--(a2.west);
  \draw (s3.east)--(a3.west);
  \draw (s4.east)--(a4.west);

  % transitions (horizontal bars) between steps
  \foreach \A/\B/\lbl in {s0/s1/{SB1（I0.0）},s1/s2/{T1: 6\,s 到},s2/s3/{T2: 6\,s 到},s3/s4/{T3: 6\,s 到}} {
    \draw[very thick] (\A) -- (\B);
    \coordinate (m) at ($(\A)!0.5!(\B)$);
    \draw[very thick] ($(m)+(-0.5,0)$) -- ($(m)+(0.5,0)$);
    \node[right] at ($(m)+(0.55,0)$) {\footnotesize \lbl};
  }

  % loop from S4 back to S1: transition T4 3s
  \draw[very thick] (s4.west) -- ++(-2.4,0) coordinate(l4);
  \draw[very thick] (l4) -- (l4|-s1.west) coordinate(l1);
  \draw[very thick,-{Latex}] (l1) -- (s1.west);
  \draw[very thick] ($(l4)!0.5!(l1)+(-0.5,0)$) -- ($(l4)!0.5!(l1)+(0.5,0)$);
  \node[left,align=right] at ($(l4)!0.5!(l1)+(-0.55,0)$) {\footnotesize T4: 3\,s 到\\（循环）};

  % stop transition: any step -> S0 when SB2, shown as note
  \node[draw,dashed,align=left,font=\footnotesize] at (7.2,0.4)
    {停止：按下 SB2(I0.1) 时\\ 复位运行标志与全部步，\\ 返回初始步 S0（全部停喷）。};
\end{tikzpicture}
\end{document}
